\documentclass[10pt,a4paper]{amsart}
\usepackage{amsmath,amssymb,amsthm,graphicx}
\usepackage{verbatim}
\newenvironment{reference}{\comment}{\endcomment}
%\newenvironment{reference}{}{}
\newenvironment{slogan}{\comment}{\endcomment}
\newenvironment{history}{\comment}{\endcomment}

% For commutative diagrams we use Xy-pic
\usepackage[all]{xy}

% We use 2cell for 2-commutative diagrams.
\xyoption{2cell}
% Theorem environments.
%
\theoremstyle{plain}
\newtheorem{theorem}[subsection]{Theorem}
\newtheorem{proposition}[subsection]{Proposition}
\newtheorem{lemma}[subsection]{Lemma}

\theoremstyle{definition}
\newtheorem{definition}[subsection]{Definition}
\newtheorem{example}[subsection]{Example}
\newtheorem{exercise}[subsection]{Exercise}
\newtheorem{situation}[subsection]{Situation}

\theoremstyle{remark}
\newtheorem{remark}[subsection]{Remark}
\newtheorem{remarks}[subsection]{Remarks}

\numberwithin{equation}{subsection}

% Macros
%
\def\lim{\mathop{\mathrm{lim}}\nolimits}
\def\colim{\mathop{\mathrm{colim}}\nolimits}
\def\Spec{\mathop{\mathrm{Spec}}}
\def\Hom{\mathop{\mathrm{Hom}}\nolimits}
\def\Ext{\mathop{\mathrm{Ext}}\nolimits}
\def\SheafHom{\mathop{\mathcal{H}\!\mathit{om}}\nolimits}
\def\SheafExt{\mathop{\mathcal{E}\!\mathit{xt}}\nolimits}
\def\Sch{\mathit{Sch}}
\def\Mor{\mathop{\mathrm{Mor}}\nolimits}
\def\Ob{\mathop{\mathrm{Ob}}\nolimits}
\def\Sh{\mathop{\mathit{Sh}}\nolimits}
\def\NL{\mathop{N\!L}\nolimits}
\def\CH{\mathop{\mathrm{CH}}\nolimits}
\def\proetale{{pro\text{-}\acute{e}tale}}
\def\etale{{\acute{e}tale}}
\def\QCoh{\mathit{QCoh}}
\def\Ker{\mathop{\mathrm{Ker}}}
\def\Im{\mathop{\mathrm{Im}}}
\def\Coker{\mathop{\mathrm{Coker}}}
\def\Coim{\mathop{\mathrm{Coim}}}

% Boxtimes
%
\DeclareMathSymbol{\boxtimes}{\mathbin}{AMSa}{"02}

%
% Macros for moduli stacks/spaces
%
\def\QCohstack{\mathcal{QC}\!\mathit{oh}}
\def\Cohstack{\mathcal{C}\!\mathit{oh}}
\def\Spacesstack{\mathcal{S}\!\mathit{paces}}
\def\Quotfunctor{\mathrm{Quot}}
\def\Hilbfunctor{\mathrm{Hilb}}
\def\Curvesstack{\mathcal{C}\!\mathit{urves}}
\def\Polarizedstack{\mathcal{P}\!\mathit{olarized}}
\def\Complexesstack{\mathcal{C}\!\mathit{omplexes}}
% \Pic is the operator that assigns to X its picard group, usage \Pic(X)
% \Picardstack_{X/B} denotes the Picard stack of X over B
% \Picardfunctor_{X/B} denotes the Picard functor of X over B
\def\Pic{\mathop{\mathrm{Pic}}\nolimits}
\def\Picardstack{\mathcal{P}\!\mathit{ic}}
\def\Picardfunctor{\mathrm{Pic}}
\def\Deformationcategory{\mathcal{D}\!\mathit{ef}}
\usepackage{cleveref}
\crefname{originalreference}{original reference}{original references}
\Crefname{originalreference}{Original reference}{Original references}
\makeatletter
\newlabel{properties-lemma-locally-constructible}{{054C}{source}{}{sourcepointer-1}{}}
\newlabel{properties-lemma-locally-constructible@cref}{{[originalreference][0][]054C}{[1][0][]source}{}{}{}}
\makeatother
\crefname{originalreference}{original reference}{original references}
\Crefname{originalreference}{Original reference}{Original references}
\makeatletter
\newlabel{morphisms-lemma-inverse-image-constructible}{{054I}{source}{}{sourcepointer-2}{}}
\newlabel{morphisms-lemma-inverse-image-constructible@cref}{{[originalreference][0][]054I}{[1][0][]source}{}{}{}}
\makeatother
\crefname{originallemma}{original lemma}{original lemmas}
\Crefname{originallemma}{Original lemma}{Original lemmas}
\makeatletter
\newlabel{lemma-points-cartesian}{{03H4}{source}{}{sourcepointer-3}{}}
\newlabel{lemma-points-cartesian@cref}{{[originallemma][0][]03H4}{[1][0][]source}{}{}{}}
\makeatother
\crefname{originalreference}{original reference}{original references}
\Crefname{originalreference}{Original reference}{Original references}
\makeatletter
\newlabel{morphisms-theorem-chevalley}{{054K}{source}{}{sourcepointer-4}{}}
\newlabel{morphisms-theorem-chevalley@cref}{{[originalreference][0][]054K}{[1][0][]source}{}{}{}}
\makeatother
\title{Locally constructible subsets of the topological space characterized on etale covers}
\author{The Stacks Project Authors}
\date{stacks-project:spaces-properties.tex:local-snapshot-20261012; GNU FDL 1.2 or later}
\begin{document}
\maketitle
\noindent\textit{Editorial note.} All delivered mathematical text is the author's own, copied verbatim from the pinned Stacks Project chapter file spaces-properties.tex; no statement or proof step was rewritten, reordered or removed.\par\medskip
\noindent\textit{Editorial note.} Span recomputation by label pairing (LANE\_BRIEF 10.2): in this chapter the author's proof environment follows the statement environment as a standalone block; the spec covers the complete environment pair (statement span, proof span) as recorded in the pool candidate.\par\medskip
\noindent\textit{Editorial note.} Occupation scan: combined claim\_key + statement\_anchor + generic-line scan over all registration files (active and withdrawn) found 13 S10-geom published cards for spaces-properties.tex (prop-occupied.json, 2026-10-14); those 13 labels are excluded from this batch by design. For the 77 free labels: 0 occupied. The registry mentions field (7557 files) and focused-seeds (645) both point to the same S10-geom cards (dual-source check), no other occupation.\par\medskip
\ifdefined\hypertarget\hypertarget{sourcepointer-1}{}\fi
\noindent\textit{Source reference.} Original reference 054C: sources/stacks-chow/tags.txt, source lines 6635--6635. Cross-chapter standard prerequisite cited by the author inside the counted span; it is used as a pointer and its own proof is not part of this item.\par\medskip
\ifdefined\hypertarget\hypertarget{sourcepointer-2}{}\fi
\noindent\textit{Source reference.} Original reference 054I: sources/stacks-chow/tags.txt, source lines 6641--6641. Cross-chapter standard prerequisite cited by the author inside the counted span; it is used as a pointer and its own proof is not part of this item.\par\medskip
\ifdefined\hypertarget\hypertarget{sourcepointer-3}{}\fi
\noindent\textit{Source reference.} Original lemma 03H4: sources/stacks-chow/tags.txt, source lines 4553--4553. Cross-chapter standard prerequisite cited by the author inside the counted span; it is used as a pointer and its own proof is not part of this item.\par\medskip
\ifdefined\hypertarget\hypertarget{sourcepointer-4}{}\fi
\noindent\textit{Source reference.} Original reference 054K: sources/stacks-chow/tags.txt, source lines 6643--6643. Cross-chapter standard prerequisite cited by the author inside the counted span; it is used as a pointer and its own proof is not part of this item.\par\medskip
\begin{lemma}
\label{lemma-locally-constructible}
Let $S$ be a scheme. Let $X$ be an algebraic space over $S$.
Let $E \subset |X|$ be a subset. The following are equivalent
\begin{enumerate}
\item for every \'etale morphism $U \to X$ where $U$ is a
scheme the inverse image of $E$ in $U$ is a locally constructible
subset of $U$,
\item for every \'etale morphism $U \to X$ where $U$ is an
affine scheme the inverse image of $E$ in $U$ is a constructible
subset of $U$,
\item for some surjective \'etale morphism $U \to X$ where $U$ is a
scheme the inverse image of $E$ in $U$ is a locally constructible
subset of $U$.
\end{enumerate}
\end{lemma}

\begin{proof}
By Properties, Lemma \ref{properties-lemma-locally-constructible}
we see that (1) and (2) are equivalent. It is immediate that (1)
implies (3). Thus we assume we have a surjective \'etale morphism
$\varphi : U \to X$ where $U$ is a scheme such that $\varphi^{-1}(E)$
is locally constructible. Let $\varphi' : U' \to X$ be another
\'etale morphism where $U'$ is a scheme. Then we have
$$
E'' = \text{pr}_1^{-1}(\varphi^{-1}(E)) = \text{pr}_2^{-1}((\varphi')^{-1}(E))
$$
where $\text{pr}_1 : U \times_X U' \to U$ and
$\text{pr}_2 : U \times_X U' \to U'$ are the projections.
By Morphisms, Lemma \ref{morphisms-lemma-inverse-image-constructible}
we see that $E''$ is locally constructible in $U \times_X U'$.
Let $W' \subset U'$ be an affine open. Since $\text{pr}_2$ is
\'etale and hence open, we can choose a quasi-compact open
$W'' \subset U \times_X U'$ with $\text{pr}_2(W'') = W'$.
Then $\text{pr}_2|_{W''} : W'' \to W'$ is quasi-compact.
We have $W' \cap (\varphi')^{-1}(E) = \text{pr}_2(E'' \cap W'')$
as $\varphi$ is surjective, see Lemma \ref{lemma-points-cartesian}.
Thus $W' \cap (\varphi')^{-1}(E) = \text{pr}_2(E'' \cap W'')$
is locally constructible by
Morphisms, Theorem \ref{morphisms-theorem-chevalley} as desired.
\end{proof}

\end{document}
